\section{Capítulo~\ref{sec:dofaalt}. Outras teorias do \nt{DOPA}}

Cada um dos quadros ampliados se apoia em suas próprias medições diretas da dopamina (dispersão dos pontos de reversão, respostas ao desagradável e ao novo, crescimento lento, respostas à troca de sabor), mas as fórmulas que as explicam são teorias, e entre duas delas os experimentos por enquanto se contradizem.

{\footnotesize
\setlength{\tabcolsep}{3pt}\renewcommand{\arraystretch}{1.15}
\begin{longtable}{>{\raggedright}p{0.5cm}>{\raggedright}p{4.0cm}>{\raggedright}p{5.6cm}>{\raggedright}p{2.3cm}>{\raggedright\arraybackslash}p{1.7cm}}
\toprule
N.º & Proposição & Experimento e o que foi medido & Fonte & Status \\
\midrule\endfirsthead
\toprule
N.º & Proposição & Experimento e o que foi medido & Fonte & Status \\
\midrule\endhead
1 & A regra assimétrica~\eqref{eq:asymmetric} converge para o ponto~\eqref{eq:expectile}; com $\kappa=1/2$ é a média; para uma gota com probabilidade $1/2$, $V=\kappa$ (fig.~\ref{fig:expectiles}) & O ponto~\eqref{eq:expectile} é um expectil, solução de um problema de mínimos quadrados com pesos diferentes para desvios positivos e negativos; para o exemplo do capítulo, a dedução está no próprio capítulo. & \cite{NeweyPowell1987}; dedução no capítulo & teoria \\
2 & Neurônios \nt{DOPA} diferentes têm pontos de reversão diferentes, ordenados pela assimetria (proposição~\ref{prop:expectile}) & Camundongos, neurônios \nt{DOPA} da área tegmental ventral marcados optogeneticamente, recompensas de tamanhos diferentes: o ponto em que o pico dá lugar à queda difere entre neurônios; nos neurônios com respostas mais assimétricas ele é mais alto; das respostas do grupo se reconstrói a forma da distribuição das recompensas. Que essa seja a função principal da dispersão é hipótese. & \cite{Dabney2020} & medido \\
3 & Parte dos neurônios \nt{DOPA} responde ao desagradável com um pico & Macacos: alguns neurônios \nt{DOPA} são excitados pelo sinal de recompensa e inibidos pelo sinal de um jato de ar no olho; outros são excitados por ambos; os grupos ficam em partes diferentes do mesencéfalo. & \cite{Matsumoto2009, BrombergMartin2010} & medido \\
4 & O módulo do erro ou a novidade define a taxa de aprendizado & Nos trabalhos citados não há experimento direto em que o sinal de importância dos neurônios \nt{DOPA} mude a taxa de aprendizado; a revisão propõe o papel de sinal ``atenção, é importante''. & \cite{BrombergMartin2010} & hipótese \\
5 & Um fio separado para o eixo do perigo & Camundongos: os neurônios \nt{DOPA} que vão à extremidade posterior do estriado respondem a estímulos novos e fortes, e não à recompensa; sua destruição enfraquece a evitação de um objeto novo. & \cite{Menegas2018} & medido \\
6 & Tempo ótimo de ação $\tau_a^*=\sqrt{C/\bar R}$~\eqref{eq:vigor} & Mínimo da soma $C/\tau_a+\bar R\tau_a$; dedução no capítulo. No modelo original, a velocidade das ações é dada pelo preço do tempo, igual à recompensa média. & \cite{Niv2007}; dedução no capítulo & teoria \\
7 & O nível lento de dopamina leva $\bar R$ e define a velocidade das ações (proposição~\ref{prop:multiplex}) & Ratos, microdiálise e voltametria no núcleo accumbens: o nível de dopamina, minuto a minuto, acompanha a frequência das recompensas e a velocidade das ações. Em humanos, a L-DOPA reforça a dependência do tempo de reação em relação à frequência média de recompensas. Que o nível lento seja exatamente $\bar R$ não foi demonstrado. & \cite{Hamid2016, Beierholm2013vigor} & hipótese \\
8 & O nível lento vem de duas fontes: a liberação muda nas saídas sem mudança na taxa de disparos, mas o crescimento lento também aparece na própria taxa & Ratos, uma mesma tarefa: a liberação no núcleo accumbens acompanha a expectativa de recompensa em mudança, mas a taxa de disparos dos neurônios \nt{DOPA} identificados da área tegmental ventral não. Camundongos num corredor virtual (linha~10): o crescimento suave ao se aproximar da recompensa aparece também nos disparos dos neurônios \nt{DOPA} identificados. & \cite{Mohebi2019, Kim2020} & medido \\
9 & A dopamina sobe suavemente enquanto o animal se aproxima de uma recompensa distante & Ratos num labirinto, voltametria no estriado: a concentração de dopamina sobe em segundos à medida que o alvo se aproxima; a inclinação depende da distância e do tamanho da recompensa. & \cite{Howe2013ramp, Hamid2016} & medido \\
10 & O crescimento é o erro $\delta$ com refinamento da estimativa~\eqref{eq:ramp}; um salto no corredor dá um pico (fig.~\ref{fig:ramp}) & A fórmula e o número $\delta=0.75\,V$ por segundo são dedução do capítulo. Experimento: camundongos num corredor virtual, transporte instantâneo para mais perto do alvo; os disparos dos corpos celulares, o cálcio dos corpos e das saídas e a concentração no estriado respondem como erro, e não como a expectativa $V$. Qual das duas explicações vale em todas as saídas está em discussão. & \cite{Kim2020} & medido \\
11 & Olhar para trás: a dopamina informa uma medida de ligação causal~\eqref{eq:retro} & Camundongos, liberação de dopamina no núcleo accumbens em experimentos em que essa teoria e o erro~\eqref{eq:ctd} preveem coisas diferentes: em vários experimentos a liberação seguiu a primeira. & \cite{Jeong2022} & hipótese \\
12 & A resposta da dopamina durante o aprendizado se desloca gradualmente da recompensa para o sinal antecipador & Camundongos, sinal das saídas dos neurônios \nt{DOPA} no estriado ventral ao longo do aprendizado: o pico passa gradualmente do instante da recompensa para o do sinal antecipador, como exige o aprendizado por~\eqref{eq:ctd}. & \cite{Amo2022} & medido \\
13 & Os neurônios \nt{DOPA} respondem à troca do sabor da recompensa com o mesmo valor & Ratos, registro de neurônios da área tegmental ventral: trocar um suco por outro de valor igual provoca resposta, embora o erro~\eqref{eq:ctd} seja zero. & \cite{Takahashi2017} & medido \\
14 & Picos artificiais ensinam a ligação entre dois estímulos neutros & Ratos, experimento de pré-condicionamento sensorial de dois estímulos sem recompensa: a ativação optogenética dos neurônios \nt{DOPA} na transição do primeiro para o segundo cria a ligação, e a desativação impede que seja aprendida. & \cite{Sharpe2017} & medido \\
15 & Vetor de erros por característica~\eqref{eq:vectortd} & Teoria do erro de predição de características; não há verificação direta da forma vetorial. & \cite{Gardner2018} & hipótese \\
16 & Neurônios \nt{DOPA} diferentes numa mesma tarefa levam grandezas diferentes & Camundongos num corredor virtual, imagem de dois fótons de neurônios \nt{DOPA}: uns ligados à recompensa, outros à velocidade e à rotação, outros aos sinais antecipadores e à precisão da escolha. & \cite{Engelhard2019} & medido \\
17 & Um feixe em vez de um fio (proposição~\ref{prop:bundle}) & Resumo das linhas~2--16: cada decomposição foi medida separadamente; não há experimento que mostre um quadro único em que todas sejam partes de um mesmo sinal. & --- & hipótese \\
\bottomrule
\end{longtable}}
